% Page layout
\usepackage{geometry}
\geometry{a4paper,left=3.18cm,right=3.18cm,top=2.54cm,bottom=2.54cm}

% Figures, tables, and colors
\usepackage{graphicx}
\usepackage[table,xcdraw]{xcolor}
\usepackage{array}
\usepackage{booktabs}
\usepackage{colortbl}
\usepackage{float}
\usepackage{footnote}
\usepackage{longtable}
\usepackage{lscape}
\usepackage{makecell}
\usepackage{multirow}
\usepackage{subfigure}
\usepackage{threeparttable}
\usepackage{caption}

\renewcommand{\figurename}{图}
\renewcommand{\tablename}{表}

% Mathematics and theorem environments
\usepackage{amsmath,amsthm,amsfonts,amssymb,mathrsfs}
\usepackage{bbding}
\usepackage[framemethod=tikz]{mdframed}

\newtheorem{theorem}{定理}[section]
\newtheorem{definition}{定义}[section]
\newtheorem{lemma}{引理}[section]
\newtheorem{proposition}{命题}[section]
\newtheorem{property}{性质}[section]
\newtheorem{axiom}{公理}[section]
\newtheorem{corollary}{推论}[section]
\newtheorem{exercise}{习题}[section]
\renewcommand{\proofname}{\bfseries 证明}
\allowdisplaybreaks[4]

% Code and algorithms
\usepackage{listings}
\usepackage{minted}
\usepackage{algorithm}
\usepackage{algorithmicx}
\usepackage{algpseudocode}

\floatname{algorithm}{Algorithm}
\renewcommand{\algorithmicrequire}{\textbf{Input:}}
\renewcommand{\algorithmicensure}{\textbf{Output:}}

\makeatletter
\newenvironment{breakablealgorithm}
{
  \begin{center}
  \refstepcounter{algorithm}
  \hrule height.8pt depth0pt \kern2pt
  \renewcommand{\caption}[2][\relax]{
    {\raggedright\textbf{算法~\thealgorithm} ##2\par}
    \ifx\relax##1\relax
      \addcontentsline{loa}{algorithm}{\protect\numberline{\thealgorithm}##2}
    \else
      \addcontentsline{loa}{algorithm}{\protect\numberline{\thealgorithm}##1}
    \fi
    \kern2pt\hrule\kern2pt
  }
}
{
  \kern2pt\hrule\relax
  \end{center}
}
\makeatother

% Chapter and contents formatting
\usepackage{titletoc}
\usepackage[explicit]{titlesec}

\titlecontents{chapter}[0pt]{\addvspace{1.5pt}\filright\bfseries}
  {\contentspush{\thecontentslabel \quad}}
  {}{\titlerule*[8pt]{.}\contentspage}

\ctexset{
  part/name = {第, 部分},
  part/number = \Roman{part}
}

% Diagrams
\usepackage{tikz}
\usepackage{pgfplots}
\pgfplotsset{compat=1.18}
\usetikzlibrary{
  arrows,
  arrows.meta,
  backgrounds,
  calc,
  positioning,
  shapes.geometric
}

% Lists, spacing, and appendices
\usepackage{enumitem}
\usepackage{setspace}
\usepackage[toc,page]{appendix}

\setenumerate[1]{itemsep=0pt,partopsep=0pt,parsep=\parskip,topsep=5pt}
\setitemize[1]{itemsep=0pt,partopsep=0pt,parsep=\parskip,topsep=5pt}
\setdescription{itemsep=0pt,partopsep=0pt,parsep=\parskip,topsep=5pt}

% References and links
\usepackage[round,sort,comma,numbers]{natbib}
\usepackage[
  colorlinks=true,
  linkcolor=blue,
  citecolor=blue,
  urlcolor=blue
]{hyperref}

\renewcommand\cite[1]{(\citeauthor{#1}, \citeyear{#1})}
\renewcommand\citep[1]{\citeyear{#1}}
\renewcommand\citet[1]{\citeauthor{#1}}
\bibliographystyle{plainnat}

% Utilities
\usepackage{framed}
\usepackage{stfloats}

\newcommand{\tabincell}[2]{\begin{tabular}{@{}#1@{}}#2\end{tabular}}

% Quote block: left accent bar + light background
\mdfdefinestyle{blockquotestyle}{%
  skipabove=\medskipamount,
  skipbelow=\medskipamount,
  leftline=true,
  rightline=false,
  topline=false,
  bottomline=false,
  linewidth=3pt,
  linecolor=black!28,
  backgroundcolor=black!4,
  innertopmargin=8pt,
  innerbottommargin=8pt,
  innerleftmargin=12pt,
  innerrightmargin=6pt,
}
\renewenvironment{quote}{%
  \begin{mdframed}[style=blockquotestyle]%
  \small\linespread{1.3}\selectfont%
}{%
  \end{mdframed}%
}

\definecolor{bg}{rgb}{0.95,0.95,0.95}
\definecolor{myBGcolor}{RGB}{204,204,255}

% Headers and footers
\usepackage{fancyhdr}
\pagestyle{fancy}
\setlength{\headheight}{13pt}
\cfoot{}
\rfoot{\thepage}
\renewcommand{\headrulewidth}{0.4pt}
\renewcommand{\footrulewidth}{0.4pt}
